% !PRINTED 97
\weberpage{97}

The transformations $S$ therefore compose according to the same rule as the linear substitutions and matrices which we considered in the sixth section of the second volume. The degree of a composed transformation is equal to the product of the degrees of the components. These matrices are here bound to the assumption that their elements are integers and their determinant is positive.

These properties are preserved under the composition of the transformations. Nevertheless the totality $\mathfrak{S}$ of the transformations $S$ forms no group, any more than the totality of the natural numbers forms a group under composition by multiplication; for when $S'$, $S''$ are given it is not always possible to determine an $S$ that satisfies condition (6), which yet (according to Bd. II, \S\ 1, 4.) would be required for a group.

By the special transformation of degree $m^2$:
\[
M = \begin{pmatrix} \pm m, & 0 \\ 0, & \pm m \end{pmatrix}
\]
the periods $\omega_1, \omega_2$ pass into $\omega_1' = \omega_1/m$, $\omega_2' = \omega_2/m$, and the period ratio $\omega = \omega_2/\omega_1$ remains unchanged. These transformations are called \gesperrt{multiplications} (similarity transformations, Bd. II, \S\ 41). Among them is contained the identical substitution
\[
1 = \begin{pmatrix} 1, & 0 \\ 0, & 1 \end{pmatrix}
\]
which leaves everything unchanged and plays the role of the unit under composition.

\gesperrt{The multiplications are permutable under composition with every transformation $S$:}

\weq{7}{SM = MS.}

If in $SM$ or $MS$ one holds the transformation $S$ fixed and lets $M$ run through the totality $\mathfrak{M}$ of the multiplications, one obtains a system $\mathfrak{M}S$, which according to Bd. II, \S\ 46 would have to be designated a collineation. If $S_1$ and $S_2$ belong to a collineation $C$, and $S_1'$ and $S_2'$ to a collineation $C'$, then $S_1'S_1$ and $S_2'S_2$ also belong to the same collineation $C''$. In this way, by setting $C'' = C'C$, one can compose the collineations. Under this composition the colli-
%CONT>>
